\documentclass{article}

\usepackage[english]{babel}
\usepackage{csquotes}
\usepackage[backend=biber, style=numeric]{biblatex}
\usepackage[colorlinks=true, allcolors=blue]{hyperref}
\addbibresource{bibliography/references.bib}   % with the extension

\title{Portable 3D Scanner for Digital Reconstruction}
\author{Raul Casquinha}
\date{\today}



\begin{document}

\maketitle

\cite{wang_mems_2020}.

\section{Abstract}

\tableofcontents

\section{Introduction}

\section{State of the Art}

\section{Methodology}

\section{Results}

\section{Discussion}

\section{Conclusion}

\section{Brainstorming}

Ficheiro de armazenamento da point cloud
Transmissão da point cloud para o computador
Interface de visualização da point cloud
Programa para reconstrução da point cloud em malha
Através de um ponto no espaço, determinar a posição do scanner (Geometria que permite a colocação do ponto no espaço)
Montagem de um protótipo handheld scanner (Esqueleto apenas com buracos para os sensores e a placa)


\printbibliography

\end{document}
